\section{Thesis outline}\label{sec:introduction:thesis_outline}
The thesis proceeds from establishing the evidential case for a missing abstraction tier, through the formal design and proof of an artefact that inhabits it.
The thesis proceeds to a demonstrated compiler that executes programs on real hardware, culminating in a synthesis that answers the central research question in full.
The thesis is structured as follows.

\begin{itemize}

    \item \textbf{Chapter~\ref{ch:literature_review} Literature Review:}
    Surveys the quantum programming language landscape; develops the abstraction framework and taxonomy used throughout the thesis; and establishes the gap statement and requirements that motivate the language design.

    \item \textbf{Chapter~\ref{ch:methodology} Methodology:}
    Establishes the Design Science Research framework governing the thesis; maps the research design to the DSR process; and defines the evaluation criteria and acceptance thresholds against which the artefacts are assessed.

    \item \textbf{Chapter~\ref{ch:pub_qpl_review} Quantum Programming Language Systematic Review:}
    Presents the first systematic review of quantum programming languages using a formal taxonomy and evaluation rubric.
    Introduces and operationalises quantum bias and establishes the evidential basis for the quantum-implicit tier as an unoccupied position in the design space.

    \item \textbf{Chapter~\ref{ch:pub_language_design} Language Specification:}
    Defines the syntax, operational semantics and type system of a quantum-implicit programming language.
    Demonstrates that representative quantum algorithms can be expressed without quantum-mechanical vocabulary.

    \item \textbf{Chapter~\ref{ch:pub_formal_theory} Type System and Formal Theory:}
    Proves type safety (progress and preservation) for the core language and semantic preservation for the compilation pipeline, establishing that the compiler faithfully preserves the meaning of source programs.

    \item \textbf{Chapter~\ref{ch:pub_compilation} Compiler:}
    Presents a compiler from the quantum-implicit language to gate-based quantum hardware and demonstrates end-to-end executability on NISQ hardware.
    This provides the practical existence proof that the quantum-implicit tier is operationally realisable.

    \item \textbf{Chapter~\ref{ch:results_evaluation} Results and Evaluation:}
    Consolidates the artefact contributions against the evaluation criteria and answers the central research question.

    \item \textbf{Chapter~\ref{ch:study_implications} Study Implications:}
    Situates the contributions in the broader landscape of quantum programming language research.

    \item \textbf{Chapter~\ref{ch:future_directions_conclusion} Future Directions and Conclusion:}
    Identifies directions for future work and closes the thesis.

\end{itemize}